% unidades/29-keigo-intro.tex
\chapter{🙇 敬語 — Introducción al Keigo}
\unidadheader{29}{敬語}{Keigo}{Hablar con respeto y humildad}

\begin{tcolorbox}[vocabbox, title={📌 Vocabulario Nuevo}]
\begin{tabularx}{\linewidth}{llX}
\toprule
\textbf{Japonés} & \textbf{Español} & \textbf{Tipo} \\
\midrule
\vocabitem{敬語}{けいご}{lenguaje respetuoso}{sust.}
\vocabitem{丁寧語}{ていねいご}{lenguaje cortés (desu/masu)}{sust.}
\vocabitem{尊敬語}{そんけいご}{lenguaje honorífico (exaltar al otro)}{sust.}
\vocabitem{謙譲語}{けんじょうご}{lenguaje humilde (bajar a uno mismo)}{sust.}
\vocabitem{お客様}{おきゃくさま}{cliente / invitado}{sust.}
\vocabitem{社長}{しゃちょう}{presidente de empresa}{sust.}
\vocabitem{部長}{ぶちょう}{jefe de departamento}{sust.}
\vocabitem{召し上がる}{めしあがる}{comer / beber (honorífico)}{v. gr.1}
\vocabitem{いらっしゃる}{—}{estar / ir / venir (honorífico)}{v. gr.1}
\vocabitem{おっしゃる}{—}{decir (honorífico)}{v. gr.1}
\bottomrule
\end{tabularx}
\end{tcolorbox}

\subsection*{🗣️ Frases Clave}

\frase{\ruby{何}{なに}を\ruby{召}{め}し\ruby{上}{あ}がりますか？}{¿Qué desea comer? (a un cliente)}
\frase{\ruby{少々}{しょうしょう}お\ruby{待}{ま}ちください。}{Espere un momento, por favor. (formal)}
\frase{\ruby{社長}{しゃちょう}はもうお\ruby{帰}{かえ}りになりました。}
  {El presidente ya se retiró a su casa.}
\frase{よろしくお\ruby{願}{ねが}い\ruby{致}{いた}します。}
  {Le ruego que me favorezca. (muy formal)}

\subsection*{⚙️ Gramática}

\patron{Honorífico (Sonkeigo): Exaltar al interlocutor}
  {お + V(stem) + になります}
  {\ej{\ruby{先生}{せんせい}は\ruby{新聞}{しんぶん}をお\ruby{読}{よ}みになります。}
    {El profesor lee el periódico.}}

\patron{Humilde (Kenjougo): Humillarse uno mismo}
  {お + V(stem) + します / いたします}
  {\ej{\ruby{荷物}{にもつ}をお\ruby{持}{も}ちします。}{Yo le cargaré su equipaje.}}

\begin{tcolorbox}[ejerciciobox, title={📝 Ejercicios Rápidos}]
\begin{enumerate}
  \item Di "Espere un momento" de tres formas: casual, polite y keigo.
  \item ¿Cuál es la diferencia entre 尊敬語 y 謙譲語?
  \item Traduce: \ruby{社長}{しゃちょう}はどちらにいらっしゃいますか？
\end{enumerate}
\vspace{0.4em}
\textbf{Respuestas:}
1) 待って / 待ってください / 少々お待ちください\quad
2) Sonkeigo = para acciones de otros (respeto); Kenjougo = para acciones propias (humildad).\quad
3) ¿Dónde se encuentra el presidente?
\end{tcolorbox}

\begin{tcolorbox}[kanjibox, title={🀄 Kanjis de la Unidad}]
\begin{tabularx}{\linewidth}{cllXX}
\toprule
\textbf{Kanji} & \textbf{On} & \textbf{Kun} & \textbf{Significado} & \textbf{Vocab} \\
\midrule
\kanjirow{敬}{けい}{うやま}{respeto}{\ruby{敬語}{けいご} / \ruby{敬意}{けいい}}
\kanjirow{様}{よう}{さま}{señor / aspecto}{\ruby{お客様}{おきゃくさま} / \ruby{様子}{ようす}}
\kanjirow{申}{しん}{もう}{decir (humilde)}{\ruby{申}{もう}す / \ruby{申請}{しんせい}}
\kanjirow{致}{ち}{いた}{hacer (humilde)}{\ruby{致}{いた}す / \ruby{一致}{いっち}}
\kanjirow{召}{しょう}{め}{llamar / comer}{\ruby{召}{め}し\ruby{上}{あ}がる}
\bottomrule
\end{tabularx}
\end{tcolorbox}
